%----------------------------------------------------------------------------------------
%	PACKAGES AND THEMES
%----------------------------------------------------------------------------------------

\documentclass[presentation,aspectratio=1610]{beamer}
\usepackage{graphicx} % Allows including images
\usepackage{booktabs} % Allows the use of \toprule,
\usepackage{caption}
\usepackage{subcaption}
\usepackage{amsmath}
\usepackage{mathtools, bm}
\usepackage{multicol}
% --------------------Chinese Support--------------------------
\usepackage{ctex}
\usepackage{xeCJK}
\usepackage{fontspec} 
\setCJKmainfont{Microsoft YaHei}
% --------------------Biber Support----------------------------
\usepackage[backend=bibtex]{biblatex}
\addbibresource{bib.bib}

\makeatletter
\let\@@magyar@captionfix\relax
\makeatother

\usefonttheme{professionalfonts}

\DeclareOptionBeamer{width}
{\PassOptionsToPackage{width=#1}{beamerouterthemesidebar}}

\DeclareOptionBeamer{hideothersubsections}
{\PassOptionsToPackage{hideothersubsections=#1}{beamerouterthemesidebar}}

\DeclareOptionBeamer{hideallsubsections}
{\PassOptionsToPackage{hideallsubsections=#1}{beamerouterthemesidebar}}

\ProcessOptionsBeamer

\mode<presentation>

\useoutertheme[height=0pt,left]{sidebar}

\usecolortheme{seahorse}
\setbeamercolor*{titlelike}{parent=structure}

\useinnertheme{rectangles}
\setbeamertemplate{frametitle}{
\leftskip=10pt\vspace{16pt}\insertframetitle
}

% ---------------Global Color Settings----------------------------
%\definecolor{whucolor}{RGB}{0 114 198}%Blue
\definecolor{whucolor}{RGB}{22 89 154}%WHUCSBlue
%\definecolor{whucolor}{RGB}{1 168 156}%Cyan
%\definecolor{whucolor}{RGB}{0 128 0}%Green
%\definecolor{whucolor}{RGB}{249 105 14}%Orange
%\definecolor{whucolor}{RGB}{102 51 153}%Purple
%\definecolor{whucolor}{RGB}{150 40 26}%Red
% ----------------------------------------------------------------

\addtobeamertemplate{sidebar left}{}{{\hspace{0.15cm}\includegraphics[width=1.3cm]{whu.png}}\vspace{0.4cm}}
\setbeamercolor{sidebar}{bg=whucolor}

\setbeamercolor{author in sidebar}{bg=white,fg=white!20}
\setbeamerfont{author in sidebar}{series=\itshape}

\setbeamerfont{title}{size = {\fontsize{26}{36}}, series=\bfseries}
\setbeamerfont{author}{size = {\fontsize{20}{40}}}
\setbeamerfont{institute}{size = {\fontsize{15}{40}}}
\setbeamerfont{date}{size = {\fontsize{18}{20}}}
\setbeamerfont{frametitle}{size = {\fontsize{20}{20}}, series=\bfseries}


\setbeamercolor{title in sidebar}{bg=white,fg=white}
\setbeamerfont{title in sidebar}{series=\bfseries}

\setbeamercolor{structure}{fg=whucolor}
\setbeamercolor{section in sidebar}{fg=white}
\setbeamercolor{subsection in sidebar}{fg=white}


\setbeamercolor{section in toc}{fg=whucolor}
\setbeamercolor{subsection in toc}{fg=whucolor}
\setbeamerfont{section in toc}{size = {\fontsize{16}{30}}}
\setbeamerfont{subsection in toc}{size = {\fontsize{12}{16}}}

\setbeamertemplate{section in toc}
{\leavevmode\leftskip=5ex%
	\llap{%
		\usebeamerfont*{section number projected}%
		\usebeamercolor{section number projected}%
		\begin{pgfpicture}{-1ex}{-0.5ex}{3ex}{4ex}
			\color{bg}
			\pgfpathcircle{\pgfpoint{0pt}{.75ex}}{1.7ex}
			\pgfusepath{fill}
			\pgftext[base]{\color{fg}\inserttocsectionnumber}
		\end{pgfpicture}\kern1.25ex%
	}%
	\inserttocsection\par}
\setbeamertemplate{subsection in toc}
{\leavevmode\leftskip=3em$\bullet$\hskip.5em\inserttocsubsection\par}


\setbeamerfont{caption name}{size = {\fontsize{12}{12}}}
\setbeamerfont{caption}{size = {\fontsize{12}{12}}}

%\usepackage[T1]{fontenc}
%\usepackage[default, osfigures]{opensans}

\usepackage{transparent}
\usebackgroundtemplate{
	% ---------------------aspectratio=1610-----------------------
	\hspace*{0.235\paperwidth}
	{\includegraphics[width=0.6\paperwidth]{whucs}}
	% ------------------------------------------------------------

	% ---------------------aspectratio=169------------------------
	%\hspace*{0.26\paperwidth}
	%{\includegraphics[width=0.55\paperwidth]{whucs}}
	% ------------------------------------------------------------

	% ---------------------aspectratio=43-------------------------
	%\hspace*{0.18\paperwidth}
	%{\includegraphics[width=0.725\paperwidth]{whucs}}
	% ------------------------------------------------------------
}

\setbeamertemplate{itemize item}[circle]
\setbeamertemplate{itemize subitem}[circle]
\setbeamertemplate{enumerate item}[circle]
\setbeamertemplate{enumerate subitem}[circle]
\setbeamercolor{block title}{bg=whucolor,fg=white}
\setbeamercolor{block title example}{fg=whucolor}

\makeatletter
\setbeamercolor{myfootlinetext}{fg=whucolor}
\newdimen\mywidth%
\setlength{\mywidth}{\paperwidth}%
\addtolength{\mywidth}{-\beamer@sidebarwidth}%
\setbeamertemplate{footline}
{
  \leavevmode%
  \hbox{%
  \begin{beamercolorbox}[wd=\beamer@sidebarwidth,ht=1ex,dp=1ex,right]{sidebar}%
  \end{beamercolorbox}
  \begin{beamercolorbox}[wd=\mywidth,ht=1ex,dp=1ex,right]{myfootlinetext}%
    \usebeamerfont{date in head/foot}\hfill\inserttitle\hspace*{2em}\insertshortauthor{}\hspace*{2em}
    \insertframenumber{} / \inserttotalframenumber\hspace*{2ex}
  \end{beamercolorbox}}%
  \vskip0pt%
}
\makeatother


\setbeamertemplate{navigation symbols}{} % To remove the navigation symbols from the bottom of all slides uncomment this line

\mode
<all>

\graphicspath{{./picture/}}
%----------------------------------------------------------------------------------------
%	TITLE PAGE
%----------------------------------------------------------------------------------------

\title[计算机学院]{仓鼠养殖与产后护理中的深度学习技术应用} % The short title appears at the bottom of every slide, the full title is only on the title page
\subtitle{————看不见我}

\author{答辩人：苍术} % Your name
\institute[WHU]{指导老师：大仓鼠}% Your instructor
\date{\today} % Date, can be changed to a custom date

\begin{document}

\begin{frame}
\titlepage % Print the title page as the first slide
\end{frame}

\begin{frame}
\frametitle{目录} % Table of contents slide, comment this block out to remove it
    \begin{columns}[T]
    	
		\begin{column}{.45\textwidth}
				\centering
				\tableofcontents[sections={1-3}]
		\end{column}
		\begin{column}{.45\textwidth}
				\centering
			\tableofcontents[sections={4-7}]
		\end{column}
	\end{columns}
\end{frame}

%----------------------------------------------------------------------------------------
%	PRESENTATION SLIDES
%----------------------------------------------------------------------------------------

%------------------------------------------------
\section{介绍} % Sections can be created in order to organize your presentation into discrete blocks, all sections and subsections are automatically printed in the table of contents as an overview of the talk
%------------------------------------------------

%\subsection{编译方式} % A subsection can be created just before a set of slides with a common theme to further break down your presentation into chunks

\begin{frame}
\frametitle{介绍}
	\begin{itemize}
	\item 改编自如下Beamer主题：\url{https://github.com/dscroft/coventry_beamer}
	\item 改造自模板：\url{https://github.com/xzsunbest/whurs-beamer}
    \item {编译方式}
	    \begin{itemize}
	    	\item 推荐安装完整版的 TeXLive
	    	\item 编译方式为：xelatex -> bibtex -> xelatex*2
	    \end{itemize}
    \item 请参考 \LaTeX 和 Beamer 用户文档 
    \item 内置七种主题颜色（CS蓝、蓝、青、绿、橙、紫、红），默认采用CS蓝
    \item 默认长宽比为16:10，提供16:9与4:3选项对应的背景水印排布方式
    \item 可在whucs.tex内进行修改
  \end{itemize}
\end{frame}

%------------------------------------------------
\section{测试章节一} % Sections can be created in order to organize your presentation into discrete blocks, all sections and subsections are automatically printed in the table of contents as an overview of the talk
%------------------------------------------------

\subsection{分块测试} % A subsection can be created just before a set of slides with a common theme to further break down your presentation into chunks

\begin{frame}
\frametitle{分块测试}
	\begin{block}{分块 1}
	这是第1分块。
	\end{block}

	\begin{block}{Block 2}
	This is the second block.
	\end{block}

	\begin{block}{Block 3}
	仓鼠是无情的生物
	\end{block}
\end{frame}

%------------------------------------------------

\subsection{分栏测试}

\begin{frame}
\frametitle{分栏测试}
\begin{columns}[c] % The "c" option specifies centered vertical alignment while the "t" option is used for top vertical alignment

\column{.45\textwidth} % Left column and width
\textbf{Heading}
\begin{enumerate}
\item Statement(陈述)
\item Explanation(解释)
\item Example(示例)
\end{enumerate}

\column{.5\textwidth} % Right column and width
Wuhan University is in Wuhan, Hubei.
It is one of the most prestigious and selective universities in China, which has been selected as a Chinese Ministry of Education Class A Double First Class University.
It was one of the four elite universities in the republican period and also one of the oldest universities in China.

\end{columns}
\end{frame}

%------------------------------------------------

\section{测试章节二}
\subsection{表格测试}

\begin{frame}
\frametitle{表格测试}
\begin{table}
\begin{tabular}{l l l}
\toprule
\textbf{Treatments} & \textbf{Response 1} & \textbf{Response 2}\\
\midrule
Treatment 1 & 0.0003262 & 0.562 \\
Treatment 2 & 0.0015681 & 0.910 \\
Treatment 3 & 0.0009271 & 0.296 \\
\bottomrule
\end{tabular}
\caption{测试表格}
\end{table}
\end{frame}

\section{测试章节三}
%------------------------------------------------

\subsection{公式测试}

\begin{frame}
\frametitle{公式测试}
Given $ g:[0,∞) \to\mathbb{R} $, With $ g(0)=0 $, derive the formula

\begin{equation}
	u(x,t)=\frac{x}{\sqrt{4\pi}}\int_{0}^{t}\frac{1}{(t-s)^{\frac{3}{2}}}e^{\frac{-x^2}{4(t-s)}}g(s)ds
\end{equation}

for a solution of the initial/boundary value problem 

\[
	\begin{cases}
		u_t-u_{xx}  = 0   & \mbox{in }  \mathbb{R}_+ \times(0,\infty)\\
		u           = 0   & \mbox{on }  \mathbb{R}_+ \times\{t=0\}\\
		u           = g   & \mbox{on }  \{x=0\}\times[0,\infty)
	\end{cases}
\]

(Hint: Let $v(x, t) := u(x, t) − g(t)$ and extend $v$ to $\{x < 0\}$ by odd reflection.)
\end{frame}

%------------------------------------------------

\section{测试章节四}
\subsection{代码测试}

\begin{frame}[fragile] % Need to use the fragile option when verbatim is used in the slide
\frametitle{代码测试}
\begin{example}[main.cpp]
\begin{verbatim}
#include<iostream>
using namespace std;
int main(){
    cout<<"Hello World!"<<endl;
    return 0;
}
\end{verbatim}
\end{example}
\end{frame}

\section{测试章节五}
%------------------------------------------------

\subsection{图片测试}

\begin{frame}
\frametitle{图片测试}
	\begin{figure}
		\begin{subfigure}{0.45\textwidth}
		  \centering
		  \includegraphics[width=0.6\linewidth]{image1.jpg}
		  \caption{1a}
		  \label{fig:sfig1}
		\end{subfigure}%
		\begin{subfigure}{0.45\textwidth}
		  \centering
		  \includegraphics[width=0.6\linewidth]{image2.jpg}
		  \caption{1b}
		  \label{fig:sfig2}
		\end{subfigure}
	\end{figure}
	\begin{figure}
	\begin{subfigure}{0.45\textwidth}
		\centering
		\includegraphics[width=0.6\linewidth]{image3.jpg}
		\caption*{1a}
		\label{fig:sfig1}
	\end{subfigure}%
	\begin{subfigure}{0.45\textwidth}
		\centering
		\includegraphics[width=0.6\linewidth]{image4.jpg}
		\caption*{1b}
		\label{fig:sfig2}
	\end{subfigure}
	\vspace{-10pt}
	\caption*{可爱的仓鼠们}
	\label{fig:fig}
	\end{figure}
\end{frame}

\section{测试章节六}
%------------------------------------------------

\subsection{引用测试}
\begin{frame}
 	\frametitle{参考文献注意事项}
 	如果你的参考文献编译失败了，可以尝试：
 	\begin{enumerate}
 		\item 清空缓存文件（可调用delete-temp-files.bat）
 		\item 考虑是否有bibtex支持
 	\end{enumerate}
\end{frame}
\begin{frame}[t,allowframebreaks]
	\frametitle{参考文献}
	\nocite{*}
	\printbibliography
\end{frame}

%------------------------------------------------

\begin{frame}{}
	\centering
		\Huge\bfseries
	\textcolor{whucolor}{The End}\\
	\textcolor{whucolor}{Q\&A}\\
\end{frame}

%----------------------------------------------------------------------------------------

\end{document}
